\section{模型评价与推广}

\subsection{模型优点}

\begin{enumerate}
  \item \textbf{物理口径统一。} 四问均在区间电量尺度上使用同一 SOC 状态方程，并显式表示弃光，避免把功率误作电量或用不等式掩盖能量去向；
  \item \textbf{严格遵守信息边界。} 日前预测只使用前序日期，日内调整冻结已执行区间，因果审计未发现未来数据泄露；
  \item \textbf{费用机制可解释。} 计划费、调整费和紧急费分别核算，能够解释预报精度提高为何不必然带来总费用下降；
  \item \textbf{结果可复核。} 全部线性规划均为凸问题，最大能量平衡残差不超过
  $8.1\times10^{-12}~\mathrm{kWh}$，SOC、充放电边界和跨日连续性均通过回代检查。
\end{enumerate}

\subsection{模型缺点}

\begin{enumerate}
  \item 预测器强调可解释性和因果安全，未引入天气、节假日等外生变量；在极端光伏和负载日仍可能发生较多紧急购电；
  \item 储能模型未考虑循环寿命、容量衰减、启停成本和充放电功率随 SOC 变化等工程细节；
  \item 终端电量以边际价格近似估值，全年最后一天的实际 SOC 仍受预测误差影响；
  \item 光伏小时预报在小时内取常值，无法描述云团快速移动引起的 10 分钟级波动。
\end{enumerate}

\subsection{推广}

该框架可推广到含风电、可控负荷、电动汽车和多类储能的园区微网。增加设备时，只需在统一能量平衡中加入相应流量、容量和成本项；若允许售电，可增加非负售电变量并采用买卖价差约束。对实际部署，可用概率预测或情景优化替代单一安全分位数，并把电池衰减成本加入目标函数，实现经济性与寿命的联合优化\cite{morales2014}。

\subsection{主要结论}

储能在固定日型下可降低 $26.90\%$ 的购电费。全年不确定环境中，因果安全裕度能够控制紧急购电风险；日内光伏更新进一步显著降低缺口，但更新次数并非越多越好。当前费用规则下，0:00 制定计划并在 6:00 更新一次是最经济的方案。实时价格预测的额外改进空间约为主方案费用的 $0.32\%$，因此优先提升光伏与负载的短期预测、优化调整触发规则，更可能获得显著收益。
